\documentclass[11pt,a4paper]{article}
\usepackage[T1]{fontenc}
\usepackage{lmodern,amsmath,amssymb,booktabs,longtable,microtype}
\usepackage[margin=25mm]{geometry}
\usepackage[hidelinks]{hyperref}
\title{Auditoria adversarial e decisões editoriais\\Pacote de submissão: grau quatro, $n=12$}
\author{Relatório técnico de revisão}
\date{16 de outubro de 2026}
\begin{document}\maketitle
\section{Parecer}
O resultado central em $n=12$ é sustentado por uma enumeração exata
assistida por condições espectrais necessárias. A revisão prepara um
manuscrito concentrado nesse resultado, com prova da validade dos filtros,
dois verificadores e checagens adicionais dos fundamentos. Não apresenta
a conjectura geral como resolvida, nem atribui à compilação Lean o alcance
de uma formalização integral.

O manuscrito amplo recebido foi substituído, na pasta de submissão, por
um texto de escopo menor e mais defensável. As afirmações gerais sobre
$n=2p^r$ e $n=4p^r$ não são usadas para sustentar a conclusão em $n=12$.
Seus testes finitos permanecem no verificador suplementar e devem ser
lidos com esse alcance. O arquivo original não foi alterado.

\section{Falhas encontradas e tratamento}
\begin{enumerate}
\item \textbf{Operador afim incorreto no texto e no Lean.}
Na convenção de bits usada, $\tau_3(9)=1$ e $\tau_3(27)=3$.
Os valores recebidos, 4 e 12, produzem
$T_4^2(u)=u\mathbin{\oplus}36$ e
$T_{12}^2(u)=u\mathbin{\oplus}45$, contradizendo os alvos 9 e 27.
Corrigimos as constantes e verificamos as órbitas e suas paridades.

\item \textbf{Uso insuficiente de \texttt{omega}.}
A afirmação $p(q-2k)+(1-2\delta)\ne0$ é verdadeira para
$p\ge3$ e $\delta\in\{0,1\}$, mas a prova recebida entregava apenas
\texttt{omega} a uma expressão com produto de variáveis. A nova prova
separa $q-2k<0$, $=0$ e $>0$, usa monotonicidade da multiplicação
por $p\ge0$ e termina por aritmética linear.

\item \textbf{Identificação artificial de uma soma em $n=8$.}
O Lean definia uma expressão $256-64k$ como se fosse a autocorrelação
real e provava $k\le2$. Essa expressão tem mínimo 128 e não demonstra
o mínimo real 32 citado no manuscrito. Retiramos essa identificação:
o teorema Lean restante prova apenas a sizígia abstrata e a contagem.
A exclusão em $n=8$ é confirmada por Walsh completo nas 1.024 funções.

\item \textbf{Ponte incompleta entre a sizígia e os geradores.}
O arquivo recebido definia $v_6$ como XOR dos outros cinco valores.
Isso prova uma consequência aritmética, mas não prova que os 43
geradores têm esses valores. Acrescentamos uma avaliação efetiva dos
43 representantes de monômios em Lean, além da terceira checagem Python.
A completude da lista de órbitas continua fora do teorema Lean.

\item \textbf{Confusão entre energia e bentness.}
O campo \texttt{bent\_all\_64\_frequencies} media 64 identidades
de energia diagonal. Renomeamos o campo para:\par
\noindent\texttt{diagonal\_energy\_all\_64\_frequencies}.\par
Uma condição necessária sem sobreviventes basta para a exclusão;
isso não torna a condição uma caracterização de bentness.

\item \textbf{Órbitas curtas em $n=2p^r$.}
A afirmação recebida de órbitas quárticas curtas de Tipos I e II para
$r\ge2$ é falsa. O estabilizador age livremente no suporte de tamanho
quatro; sua ordem divide quatro e $2p^r$, portanto só pode ser 1 ou 2.
As únicas órbitas curtas têm tamanho $p^r$ e são bi-antipodais.
O novo manuscrito contém esse argumento correto.

\item \textbf{Avaliação diagonal indevidamente universal.}
A frase $g_0(\mathbf1_6)=0$ para toda função em $n=12$ é falsa:
a órbita $P_{12,4}$ dá $g_0(u)=u_0u_3+u_1u_4+u_2u_5$, cujo valor
em $\mathbf1_6$ é 1. Essa justificativa foi eliminada.

\item \textbf{Filtros sem demonstração completa de necessidade.}
O manuscrito novo demonstra o transporte com sinal, Poisson,
as congruências de energia módulo 384 e 192, a autocorrelação
módulo 3 e 192 e a dualidade diagonal módulo 48.
Também explica a cobertura exata por uma base triangular de posto 43.
O código agora verifica a identidade da mudança de base e os pivôs,
o fechamento de todos os subespaços e cada contador final da cascata.

\item \textbf{Bibliografia e prioridade.}
O artigo relevante de Sun, Shi, Liu e Fu de 2026 estava ausente;
a referência de 2026 recebida era sobre construções 2-RS.
A referência atribuída a Sun et al. em 2018 também misturava dados.
O manuscrito de submissão usa uma bibliografia menor, pertinente,
e inclui o DOI correto: \url{https://doi.org/10.1007/s10623-026-01848-4}.
Retiramos qualificações automáticas como ``primeiro'' e ``novo demonstrado''.
\end{enumerate}

\section{Resultados e limites da reprodução}
Os postos sucessivos são
$12,18,24,27,31,34,36,37,41,41,42,43$.
Os contadores de referência são
\[
712\to361\to20158\to212296\to8688\to26496\to94936
\]
\[
\to149448\to106552\to349680\to23636\to128\to0.
\]
O crescimento de alguns contadores é esperado: cada extensão libera
novas coordenadas. Eles contam classes em quocientes diferentes,
não uma sequência de subconjuntos do mesmo tamanho inicial.

Os logs finais acompanham o pacote. O verificador principal conserva
o algoritmo recebido com correções de portabilidade, nomenclatura e
checagem de subespaços. A segunda implementação calcula as fibras
diretamente, sem transporte de Walsh, e verifica a transformação da base.
Ambas usam as mesmas condições matemáticas; portanto, sua concordância
não substitui a demonstração de necessidade dos filtros.

A terceira implementação reconstrói os geradores sem importar os dois
verificadores. Confere todas as 495 máscaras de grau quatro, as 43 órbitas,
a redução de grau e a invariância nas 64 fibras, 1.056.768 identidades
de transporte sobre os geradores, a sizígia complementar, os operadores
$T_9,T_{27}$, as imagens de postos 9, 10 e 5, e os oito setores.
Também testa Walsh completo em $n=4,6,8$.

\paragraph{Controles positivos.}
Os dezesseis exemplos quadráticos RS bent em $n=12$ passam nos
24 testes de Poisson. Esses controles não são apresentados como
validação das restrições específicas de homogeneidade quártica.
As cinquenta amostras quárticas com semente fixa verificam identidades,
não a exaustividade da busca.

\paragraph{Certificação formal.}
O arquivo revisado compila com Lean 4.15.0, sem \texttt{sorry}.
Verifica identidades auxiliares e a sizígia nos geradores explicitamente
listados. A prova de inexistência em todo o espaço $2^{43}$ continua
dependendo do argumento matemático e da execução Python/NumPy.
O log vazio de compilação é acompanhado do registro de versão e status.

\paragraph{Caso $n=10$.}
O verificador suplementar reproduz a falha do argumento módulo cinco:
$8\equiv-32$ e $-8\equiv32\pmod5$.
Reproduz também a cascata de derivadas diagonais que elimina o ramo
$|A_{31}|=8$. Os 1.625 exemplos que satisfazem todos os testes diagonais
de ordem zero mostram que \emph{esses testes específicos} não bastam.
Isso não autoriza afirmar que nenhuma propriedade possível do vetor
completo $(A_z)_z$ poderia distinguir as candidatas.

\paragraph{Dimensões maiores.}
Verificar os 245 geradores em $n=20$ permite estender por linearidade
as identidades lineares efetivamente verificadas, como cancelamento
de coeficientes e invariância. Não enumera $2^{245}$ funções, não prova
inexistência bent em $n=20$ e não é uma verificação exaustiva de uma
família infinita. A divisibilidade espectral exige ainda a passagem
matemática por grau, quociente e teorema de divisibilidade apropriado.

\section{Avaliação editorial}
O ponto forte do pacote é uma exclusão finita reproduzível em $n=12$
com todas as etapas expostas. O caso $n=4$ serve como controle elementar;
acrescentá-lo não aumenta substancialmente a novidade do resultado.
A relevância para publicação depende principalmente de a certificação
em $n=12$ e o método oferecerem algo não coberto na literatura.

O Teorema 2 de Sun et al. já cobre a condição de paridade de órbitas
curtas em $n=2p^r$ para $p\ge5$. Em $n=12$, a hipótese de tamanho do
primo daquela afirmação não se satisfaz. Isso distingue esse teorema
específico, mas não estabelece prioridade contra toda a literatura.
O pacote não promete aceitação editorial nem substitui revisão por pares.

Para submissão, o autor deve completar apenas dados editoriais que não
constam do material: periódico de destino, afiliação/endereço de contato
e declarações exigidas pelo periódico. Não foram inventadas declarações
de ineditismo, ausência de submissão simultânea ou conflitos de interesse.
Nenhum arquivo foi enviado a periódico ou repositório público nesta etapa.
\end{document}
